% =====================================================================
%  preamble.tex
%  Preamble khusus Modul Perkuliahan Desain Web.
%
%  - Bahasa Indonesia (nama bab, gambar, tabel, dll. diterjemahkan)
%  - Palet biru/teal yang terkendali (restrained blue/teal)
%  - Lingkungan pedagogis yang dapat digunakan ulang:
%      capaian, tujuanpraktik, langkah, catatan, tip,
%      peringatan, latihan, tangkapanlayar (placeholder non-gambar)
%
%  File ini di-\input dari main.tex. Jangan kompilasi sendiri.
% =====================================================================

% ---------------------------------------------------------------------
% 1. BAHASA INDONESIA
%    ElegantBook memuat lang=en; kita terjemahkan label-label inti.
% ---------------------------------------------------------------------
\renewcommand{\chaptername}{Bab}
\renewcommand{\contentsname}{Daftar Isi}
\renewcommand{\listfigurename}{Daftar Gambar}
\renewcommand{\listtablename}{Daftar Tabel}
\renewcommand{\figurename}{Gambar}
\renewcommand{\tablename}{Tabel}
\renewcommand{\bibname}{Daftar Pustaka}
\renewcommand{\indexname}{Indeks}
\renewcommand{\appendixname}{Lampiran}
\renewcommand{\proofname}{Bukti}
\renewcommand{\notename}{Catatan}
\renewcommand{\examplename}{Contoh}
\renewcommand{\exercisename}{Latihan}
\renewcommand{\problemname}{Soal}
\renewcommand{\solutionname}{Penyelesaian}
\renewcommand{\remarkname}{Catatan}
\renewcommand{\assumptionname}{Asumsi}
\renewcommand{\conclusionname}{Kesimpulan}
\renewcommand{\propertyname}{Sifat}
\renewcommand{\introductionname}{Pendahuluan}

% ---------------------------------------------------------------------
% 2. PALET WARNA — biru/teal yang terkendali
%    ElegantBook color=blue memberi structurecolor RGB(60,113,183).
%    Kita pertajam ke biru/teal yang lebih tenang tanpa mengubah
%    arsitektur warna template (main/second/third tetap dipakai).
% ---------------------------------------------------------------------
\definecolor{structurecolor}{RGB}{38, 92, 148}   % biru utama (teal-leaning)
\definecolor{main}{RGB}{38, 92, 148}             % biru utama
\definecolor{second}{RGB}{0, 128, 128}           % teal (aksen sekunder)
\definecolor{third}{RGB}{0, 105, 148}            % biru tua (aksen tersier)

% Warna bantu untuk lingkungan pedagogis (tetap dalam keluarga biru/teal)
\definecolor{pedagogyBlue}{RGB}{38, 92, 148}
\definecolor{pedagogyTeal}{RGB}{0, 128, 128}
\definecolor{pedagogyLight}{RGB}{232, 240, 248}
\definecolor{pedagogyTealLight}{RGB}{224, 240, 240}
\definecolor{pedagogyAmber}{RGB}{176, 130, 20}   % hanya untuk peringatan (kontras)
\definecolor{pedagogyAmberLight}{RGB}{250, 244, 224}

% ---------------------------------------------------------------------
% 3. LINGKUNGAN PEDAGOGIS
%    Dibangun di atas tcolorbox (sudah dimuat ElegantBook) agar
%    konsisten dengan gaya visual template.
% ---------------------------------------------------------------------

% --- 3.1 capaian (Capaian Pembelajaran) ---
\newtcolorbox{capaian}{
  enhanced,
  breakable,
  colback=pedagogyLight,
  colframe=pedagogyBlue,
  coltitle=white,
  colbacktitle=pedagogyBlue,
  fonttitle=\bfseries,
  title={Capaian Pembelajaran},
  boxrule=0.8pt,
  arc=2pt,
  left=8pt, right=8pt, top=6pt, bottom=6pt,
  before skip=10pt, after skip=10pt
}

% --- 3.2 tujuanpraktik (Tujuan Praktik) ---
\newtcolorbox{tujuanpraktik}{
  enhanced,
  breakable,
  colback=pedagogyTealLight,
  colframe=pedagogyTeal,
  coltitle=white,
  colbacktitle=pedagogyTeal,
  fonttitle=\bfseries,
  title={Tujuan Praktik},
  boxrule=0.8pt,
  arc=2pt,
  left=8pt, right=8pt, top=6pt, bottom=6pt,
  before skip=10pt, after skip=10pt
}

% --- 3.3 langkah (Langkah Kerja) — daftar bernomor otomatis ---
\newenvironment{langkah}{%
  \par\noindent\textbf{\color{pedagogyBlue}Langkah Kerja}\par\nobreak
  \begin{enumerate}[label=\textbf{\arabic*.}, leftmargin=2em, itemsep=2pt]
}{%
  \end{enumerate}
}

% --- 3.4 catatan (Catatan) ---
\newtcolorbox{catatan}{
  enhanced,
  breakable,
  colback=pedagogyLight,
  colframe=pedagogyBlue,
  coltitle=white,
  colbacktitle=pedagogyBlue,
  fonttitle=\bfseries,
  title={Catatan},
  boxrule=0.8pt,
  arc=2pt,
  left=8pt, right=8pt, top=6pt, bottom=6pt,
  before skip=10pt, after skip=10pt
}

% --- 3.5 tip (Tips) ---
\newtcolorbox{tip}{
  enhanced,
  breakable,
  colback=pedagogyTealLight,
  colframe=pedagogyTeal,
  coltitle=white,
  colbacktitle=pedagogyTeal,
  fonttitle=\bfseries,
  title={Tips},
  boxrule=0.8pt,
  arc=2pt,
  left=8pt, right=8pt, top=6pt, bottom=6pt,
  before skip=10pt, after skip=10pt
}

% --- 3.6 peringatan (Peringatan) ---
\newtcolorbox{peringatan}{
  enhanced,
  breakable,
  colback=pedagogyAmberLight,
  colframe=pedagogyAmber,
  coltitle=white,
  colbacktitle=pedagogyAmber,
  fonttitle=\bfseries,
  title={Peringatan},
  boxrule=0.8pt,
  arc=2pt,
  left=8pt, right=8pt, top=6pt, bottom=6pt,
  before skip=10pt, after skip=10pt
}

% --- 3.7 latihan (Latihan) — wadah konten bebas (prosa, daftar, kode)
%     Digunakan bab sebagai kotak latihan; isi boleh berupa paragraf,
%     itemize/enumerate, atau lstlisting. Nomor latihan: bab.nomor.
\newcounter{latihan}[chapter]
\renewcommand{\thelatihan}{\thechapter.\arabic{latihan}}
\newenvironment{latihan}[1][]{%
  \refstepcounter{latihan}%
  \begin{tcolorbox}[
    enhanced,
    breakable,
    colback=pedagogyLight,
    colframe=pedagogyBlue,
    coltitle=white,
    colbacktitle=pedagogyBlue,
    fonttitle=\bfseries,
    title={Latihan \thelatihan~#1},
    boxrule=0.8pt,
    arc=2pt,
    left=8pt, right=8pt, top=6pt, bottom=6pt,
    before skip=10pt, after skip=10pt
  ]
}{%
  \end{tcolorbox}
}

% --- 3.8 tangkapanlayar (placeholder non-gambar)
%     Input:  #1 = jalur aset yang diharapkan (mis. assets/images/bab1/halaman-utama.png)
%             #2 = keterangan (caption)
%             #3 = instruksi untuk mahasiswa
%     Menampilkan kotak placeholder yang jelas, BUKAN gambar, sehingga
%     penulis tahu aset apa yang harus disediakan.
\NewDocumentEnvironment{tangkapanlayar}{ m m m }{%
  \begin{tcolorbox}[
    enhanced,
    breakable,
    colback=pedagogyLight,
    colframe=pedagogyBlue,
    coltitle=white,
    colbacktitle=pedagogyBlue,
    fonttitle=\bfseries,
    title={Tangkapan Layar (placeholder)},
    boxrule=0.8pt,
    arc=2pt,
    left=8pt, right=8pt, top=6pt, bottom=6pt,
    before skip=10pt, after skip=10pt
  ]
  \textbf{Jalur aset yang diharapkan:} \path{#1}\par
  \textbf{Keterangan:} #2\par
  \textbf{Instruksi mahasiswa:} #3\par
  \medskip
  \noindent\fbox{%
    \parbox{0.92\linewidth}{%
      \centering
      \vspace{1.2em}
      {\large\color{pedagogyBlue}\textbf{[ Tempatkan tangkapan layar di sini ]}}\\[0.4em]
      {\small\color{black!60}Ganti blok ini dengan \texttt{\textbackslash includegraphics} setelah aset \texttt{#1} tersedia.}
      \vspace{1.2em}
    }%
  }
}{%
  \end{tcolorbox}
}

% ---------------------------------------------------------------------
% 3.9 Definisi bahasa untuk listings (css & javascript)
%     Listings di TeX Live ini tidak memuat bahasa css/javascript secara
%     bawaan; didefinisikan di sini agar blok kode bab dapat dikompilasi.
% ---------------------------------------------------------------------
\lstdefinelanguage{css}{
  morekeywords={color,background,background-color,background-image,border,
    border-radius,box-shadow,display,flex,flex-direction,flex-wrap,justify-content,
    align-items,gap,grid,grid-template-columns,grid-template-rows,grid-gap,
    margin,padding,width,height,max-width,min-width,font,font-family,font-size,
    font-weight,text-align,text-decoration,line-height,position,top,right,bottom,
    left,z-index,overflow,opacity,transition,transform,animation,media,float,
    clear,list-style,content,cursor,visibility,vertical-align,white-space},
  sensitive=true,
  morecomment=[l]{//},
  morecomment=[s]{/*}{*/},
  morestring=[b]',
  morestring=[b]"
}

\lstdefinelanguage{JavaScript}{
  keywords={break,case,catch,class,const,continue,debugger,default,delete,do,
    else,export,extends,finally,for,function,if,import,in,instanceof,let,new,
    return,super,switch,this,throw,try,typeof,var,void,while,with,yield,async,
    await,of,static,get,set,true,false,null,undefined},
  sensitive=true,
  morecomment=[l]{//},
  morecomment=[s]{/*}{*/},
  morestring=[b]',
  morestring=[b]"
}

% ---------------------------------------------------------------------
% 4. KONVENSI ASET
%    Jalur relatif terhadap folder book/.
%    - Gambar:  assets/images/<bab>/<nama>.png
%    - Kode:    assets/code/<bab>/<nama>.html
%    Gunakan \figref{label} (disediakan ElegantBook) untuk referensi.
% ---------------------------------------------------------------------
\graphicspath{{assets/images/}}

% ---------------------------------------------------------------------
% 5. HIPERLINK — warna tautan mengikuti palet biru/teal
% ---------------------------------------------------------------------
\hypersetup{
  colorlinks=true,
  linkcolor=pedagogyBlue,
  citecolor=pedagogyTeal,
  urlcolor=pedagogyTeal
}
